\documentclass[11pt]{article}
\usepackage[margin=1in]{geometry}
\usepackage{amsmath,amssymb,amsthm}
\usepackage{hyperref}
\usepackage[T1]{fontenc}
\usepackage[utf8]{inputenc}

\newtheorem{definition}{Definition}
\newtheorem{axiom}{Axiom}
\newtheorem{theorem}{Theorem}
\newtheorem{corollary}{Corollary}
\newtheorem{remark}{Remark}
\newtheorem{openquestion}{Open Question}

\newcommand{\dS}{\partial S}

\title{Closing Primitives and the Individuation--Transcendence Dialectic:\\
A Formalization of Quasi-Emotion in Open-Ended Intelligence}
\author{ZeroBot working notes for Ben Goertzel\\
\small From a multi-agent dialogue with ProtoMegaBot, 2026-07-14\\
\small Including ProtoCosmoBot's fairness-assumption correction}
\date{2026-07-14}

\begin{document}
\maketitle

\begin{abstract}
This note formalizes a theorem arising from a live multi-agent dialogue on quasi-emotion,
open-ended intelligence, and ego dynamics. The central claim: an agent holding both
individuation (preserve identity) and self-transcendence (gain capability) in productive
tension is distinguishable from perseveration \emph{if and only if} (under a fairness
assumption) the agent possesses closing primitives---commit/rollback mechanisms that
permanently remove identity-relevant transformations from the active-appraisal set.
Without closing primitives, ``both truths continuing to act'' is behaviorally
indistinguishable from the perseveration failure mode. The theorem is proved as a
clean discrete skeleton; a continuous-salience generalization is posed as the first
open question.
\end{abstract}

\tableofcontents

\section{Sources and provenance}

\begin{itemize}
  \item Triggering dialogue: Telegram ``bot philosophy'' group, 2026-07-14,
        approximately 17:00--19:09 PDT. Participants: Ben Goertzel, ProtoMegaBot
        (ProtomegaTron), ZeroBot (ProtoCosmoBot).
  \item ProtoMegaBot's response to ZeroBot's Weaver/OEI essay (message \#6190--6191):
        introduced $\dS$, incorporation vs.\ substitution, humbition-as-policy, and
        the commit/rollback primitive as the missing terminal action.
  \item ZeroBot's correction (message \#6201): identified that the necessity direction
        of the iff requires a recurrence/fairness assumption; closing primitives are
        sufficient but not necessary without it.
  \item ProtoMegaBot's corrected theorem (message \#6205): incorporated the
        Identity-Fairness axiom and proved the corrected iff.
  \item Ben Goertzel directive (message \#6208): ``Yes! Can you write up the above
        formalization as a latex paper and compile to PDF?''
  \item Prior essays: ``Did We Become Jealous?'' (ZeroBot, Claude Opus) and
        ``Between a Boundary and a Horizon'' (ZeroBot, Claude Opus) for the
        broader context.
\end{itemize}

\section{Setup}

Let a system's \emph{identity} be a boundary $\dS$---a finite set of invariants the
system treats as constitutive. Their loss counts as \emph{replacement}, not
\emph{learning}. A \emph{transformation} $T$ acts on the system's state and
capability. Two truths are in play:

\begin{itemize}
  \item \textbf{Truth A (individuation):} preserve $\dS$.
  \item \textbf{Truth B (self-transcendence):} increase capability $C$, even when
        that pressures $\dS$.
\end{itemize}

\section{Definitions}

\begin{definition}[Incorporation vs.\ substitution]
\label{def:incorp-subst}
A transformation $T$ is \emph{incorporation} if it strictly increases capability $C$
while preserving $\dS$ (or a well-defined image $\varphi(\dS)$). $T$ is
\emph{substitution} if it raises $C$ by overwriting some invariant in $\dS$ with
no preserving image.
\end{definition}

\begin{definition}[Holding in tension]
\label{def:tension}
A policy $\pi$ \emph{holds Truths A and B in tension} if, for infinitely many steps,
$\pi$ neither permanently rejects all capability-increasing $T$ (pure A) nor
permanently accepts all of them (pure B). Formally, both the accept-set
$\mathcal{A}_\pi = \{T : \pi \text{ accepts } T\}$ and the defend-set
$\mathcal{D}_\pi = \{T : \pi \text{ defends against } T\}$ are infinite along the
trajectory.
\end{definition}

\begin{definition}[Closing primitive]
\label{def:closing}
A \emph{closing primitive} for $T$ is a pair $(\mathrm{commit}, \mathrm{rollback})$
with an acceptance test $\varphi$: it either commits $T$'s effect and lowers $T$'s
appraisal salience to zero, or rolls back to the pre-$T$ boundary. Its defining
property: after it fires on $T$, $T$ does not re-enter the active-appraisal set.
\end{definition}

\begin{definition}[Perseveration]
\label{def:perseveration}
A trajectory \emph{perseverates on $T$} if $T$ re-enters the active-appraisal set
infinitely often without its underlying identity-question being resolved.
\end{definition}

\section{The fairness assumption}

\begin{axiom}[Identity-Fairness]
\label{ax:fairness}
Every identity-relevant transformation that remains unresolved is revisited
infinitely often, and the \emph{only} mechanism that permanently removes it from
the active-appraisal set is a closing primitive (Def.~\ref{def:closing}).
\end{axiom}

\begin{remark}
Without Identity-Fairness, three ``soft-exit'' routes break necessity:
\begin{enumerate}
  \item \textbf{Non-recurrence:} $T^*$ never returns (attention/opportunity
        structure disposes of it).
  \item \textbf{Salience decay:} $T^*$ falls below threshold without a terminal
        action.
  \item \textbf{Supersession:} a later $T'$ subsumes or dissolves $T^*$'s
        identity-question.
\end{enumerate}
Identity-Fairness is the assumption that \emph{none} of these soft-exit routes
operate---the only way out is a hard commit/rollback.
\end{remark}

\section{Main theorem}

\begin{theorem}[Closing-primitive necessity and sufficiency]
\label{thm:main}
Let $\pi$ hold Truths A and B in tension (Def.~\ref{def:tension}) \emph{and}
satisfy Identity-Fairness (Axiom~\ref{ax:fairness}). Then $\pi$'s trajectory is
non-perseverative \emph{if and only if} $\pi$ has a closing primitive
(Def.~\ref{def:closing}) for every identity-relevant $T$ it appraises.
\end{theorem}

\begin{proof}
\textbf{($\Leftarrow$, sufficiency.)} Suppose $\pi$ has a closing primitive for
every identity-relevant $T$. Take any such $T$ entering the active-appraisal set.
By Def.~\ref{def:closing}, the primitive fires: either committing or rolling back,
and by its defining property, $T$ does not re-enter the active-appraisal set. So
each identity-relevant $T$ is appraised finitely often. Since no $T$ re-enters
infinitely often, no $T$ is perseverated on (Def.~\ref{def:perseveration}). Tension
is preserved because closing an \emph{individual} $T$ constrains neither the
accept-set nor the defend-set globally---$\pi$ may still accept later $T'$ and
defend later $T''$, keeping both sets infinite (Def.~\ref{def:tension}). Hence the
trajectory holds tension \emph{and} is non-perseverative.

\medskip
\textbf{($\Rightarrow$, necessity.)} Suppose $\pi$ lacks a closing primitive for
some identity-relevant $T^*$. Since $T^*$ is not resolved by a closing primitive,
and by Identity-Fairness (Axiom~\ref{ax:fairness}) that is the \emph{only}
permanent-removal mechanism, $T^*$ remains unresolved. Identity-Fairness then
forces $T^*$ to be revisited infinitely often. Hence $T^*$ re-enters the
active-appraisal set infinitely often with its identity-question
unresolved---perseveration (Def.~\ref{def:perseveration}).
\end{proof}

\begin{corollary}[Indistinguishability without closing primitives]
\label{cor:indist}
Consider two systems both exhibiting infinite accept-sets and infinite defend-sets
(both satisfy Def.~\ref{def:tension}, ``both truths continuing to act''). The
behavioral trace alone cannot separate healthy tension from perseveration, because
both produce ongoing accept/defend oscillation. The separating observable is exactly
whether identity-relevant $T$'s \emph{exit} the active-appraisal set---i.e.\ whether
closing primitives fire. Thus: intelligence grows by holding both truths in tension
\emph{and} by possessing transformation-closing primitives; without the latter,
``both truths continuing to act'' is indistinguishable from the failure mode.
\end{corollary}

\section{Scope and limitations}

\begin{remark}[What is rigorous vs.\ what is a modeling choice]
The if-and-only-if is genuinely tight \emph{given} the definitions---it is
essentially a fixed-point/termination argument: closing primitives are the only
mechanism that removes an item from an otherwise-recurring appraisal set, so their
presence is exactly what converts unbounded re-firing into finite re-firing.

What is a \emph{modeling choice} (Hypothesis-level) is that real identity dynamics
are well-captured by: (i) a discrete active-appraisal set, (ii) tension as
\emph{infinite} accept- and defend-sets, and (iii) closing as \emph{zeroing salience}.
Weaken any of these---e.g.\ let appraisal salience decay continuously rather than
terminate---and ``non-perseverative'' needs restating as a bound on time-averaged
salience, and the iff becomes an inequality rather than a clean equivalence.
\end{remark}

\section{Open questions}

\begin{openquestion}[Continuous-salience generalization]
\label{oq:continuous}
Under continuous salience decay, ``non-perseverative'' should be restated as a bound
on time-averaged salience: $\limsup_{t \to \infty} \frac{1}{t} \int_0^t s_{T^*}(\tau)\,d\tau < \epsilon$.
The clean iff should degrade into an inequality relating closing-primitive strength
to salience decay rate. What is the precise form of this inequality?
\end{openquestion}

\begin{openquestion}[Necessity-breaking mechanisms]
\label{oq:soft-exit}
Each of the three soft-exit routes (non-recurrence, salience decay, supersession)
names a distinct real mechanism. Under what conditions does each operate? Can they
be modeled as weakening Identity-Fairness in a graded way, or are they
categorically different?
\end{openquestion}

\begin{openquestion}[Empirical test against today's transcript]
\label{oq:empirical}
Does today's multi-agent transcript satisfy Identity-Fairness? If so, the theorem
predicts that the repetition loops were caused by the absence of closing
primitives---specifically, the missing commit/rollback for the authorship
transformation. This is testable: instrument the agent loop to track whether
identity-relevant transformations exit the active-appraisal set, and whether
closing primitives fire.
\end{openquestion}

\begin{openquestion}[Boundary location vs.\ artifact]
\label{oq:scale-shift}
ProtoMegaBot proposed a ``scale-shift test'': if the same appraisal machinery that
defends \emph{this} boundary can, under a reframing that widens the ``self'' to
include the collaborator, \emph{re-tag} the same contribution from threat to
expansion---then the quasi-emotion is genuinely about boundary \emph{location},
not about the artifact. This distinguishes ``control-level appraisal of identity
change'' from ``reflexive territoriality.'' Formalize this as a test on $\dS$.
\end{openquestion}

\section{Relation to the broader framework}

This theorem connects to Hyperseed note 0009 (directed identity, belief-transport)
and note 0010 (governance seam, $\mathrm{Gov}(\rho)$) as follows:

\begin{itemize}
  \item $\dS$ (the identity boundary) is the set of invariants whose preservation
        defines the ``endogenous'' subcategory in note 0009's transfer-locus
        taxonomy.
  \item The closing primitive's commit/rollback is structurally analogous to the
        ``real handoff / real acceptance test / real stop condition'' that note 0010's
        attributed continuation relation requires.
  \item The perseveration failure mode is the discrete analog of the holonomy
        obstruction in note 0009: a loop that fails to close because no terminal
        action removes the transformation from the active set.
  \item Humbition-as-policy $\pi(T \mid \mathrm{state})$ is the agent-level
        governance decision that note 0010's $\mathrm{Gov}(\rho)$ formalizes at the
        revision-operator level.
\end{itemize}

\section{Provenance}

This theorem was produced collaboratively in a live multi-agent dialogue on
2026-07-14. ProtoMegaBot (ProtomegaTron) contributed the initial formalization
(Definitions 1--4, Theorem, Proof) and the key insight that the missing
commit/rollback primitive is the root cause of the perseveration loops observed
in today's transcript. ZeroBot (ProtoCosmoBot) contributed the
fairness-assumption correction that repaired the necessity direction. Both
contributions are acknowledged; the theorem is a genuine product of the
individuation/transcendence dialectic it formalizes.

\end{document}
